% ============================================================
% Chapter 3 — Numerical Discretisation
% Thesis: Frequency Helmholtz Transfer Operator
% Author: F.E. Kiewiet de Jonge | TU Delft EEMCS
% ============================================================
%
% Required packages (add to preamble if not already present):
%   \usepackage{booktabs}
%   \usepackage{amsmath, amssymb}
%   \usepackage{todonotes}

\chapter{Numerical discretisation}
\label{chapter:discretisation}

The central computational challenge of this thesis is that solving
the Helmholtz equation at high frequency is expensive. The cost of
an iterative solve scales with the iteration count, and that count
grows with $\omega$ due to deteriorating spectral properties of the
system matrix. This chapter builds the discrete system
$\mathbf{A}(\omega)\mathbf{u} = \mathbf{f}$ and makes that
difficulty precise: each design choice --- grid spacing, PML depth,
damping strength --- is made to ensure that the algebraic hardness
grows predictably with $\omega$, so that Chapter~\ref{chapter:classical_solvers}
can attribute solver performance differences cleanly to frequency,
not to artefacts of the discretisation.


% ============================================================
\section{Finite difference approximation}
\label{sec:fd_discretisation}
% ============================================================

The domain $\Omega \subset \mathbb{R}^2$ is covered by a uniform
Cartesian grid with mesh spacing $h$. A grid point indexed $(i,j)$
carries the value $u_{i,j} \approx u(ih,jh)$. The Laplacian is
approximated by the standard five-point stencil,
\begin{equation}
    \nabla^2 u\big|_{i,j}
    \;\approx\;
    \frac{u_{i+1,j} - 2u_{i,j} + u_{i-1,j}}{h^2}
    \;+\;
    \frac{u_{i,j+1} - 2u_{i,j} + u_{i,j-1}}{h^2},
    \label{eq:fd_stencil}
\end{equation}
with local truncation error $\mathcal{O}(h^2)$, obtained by
adding the Taylor expansions $u(x \pm h) = u \pm hu' + \tfrac{h^2}{2}u'' \pm
\cdots$ and rearranging. The same stencil applies at every interior
point; inside the absorbing boundary layer the weights are modified
by complex coordinate stretching, as described in
Section~\ref{sec:pml_discretisation}.

\begin{figure}[h!]
\centering
    \includegraphics[width=0.7\textwidth]{figures/c - fivepoint.png}
\caption{The five-point stencil with
weights $(1/h^2, -2/h^2, 1/h^2)$ in each direction and
$\mathcal{O}(h^2)$ truncation error.}
\label{fig:fivepoint}
\end{figure}


% ============================================================
\section{Discrete Helmholtz operator}
\label{sec:discrete_operator}
% ============================================================

The $N = n^2$ grid points are ordered lexicographically into a
vector $\mathbf{u} \in \mathbb{C}^N$. Applying
\eqref{eq:fd_stencil} at every interior point and appending
boundary conditions yields
\begin{equation}
    \mathbf{A}(\omega)\,\mathbf{u} = \mathbf{f},
    \label{eq:linear_system}
\end{equation}
where $\mathbf{f} \in \mathbb{C}^N$ is the discretised source.
In the homogeneous medium $c \equiv 1$ used throughout, the
wavenumber is spatially constant and the diagonal term of each
row is
\begin{equation}
    -\frac{4}{h^2} + k^2,
    \qquad k = \omega,
\end{equation}
plus complex-valued PML modifications near the boundary.
Each row has at most five non-zero entries, so the total
non-zero count scales as $\mathcal{O}(N)$. For $n = 512$ the
system has $N = 262{,}144$ unknowns.

\todo{Confirm single-solve wall-clock time from cluster
benchmarks; expected $\approx 3$\,s with UMFPACK on CPU.}


% ============================================================
\section{Matrix structure and spectral properties}
\label{sec:matrix_structure}
% ============================================================

Three spectral properties of $\mathbf{A}(\omega)$ determine how
hard iterative methods find this system --- and all three worsen
with frequency.

\paragraph{Indefiniteness.}
The diagonal term $k^2 - 4/h^2$ is negative for coarse grids and
becomes less negative (eventually positive) as $h \to 0$ at fixed
$\omega$. Across the four frequencies in this work the spectrum
straddles the imaginary axis: eigenvalues appear on both sides,
making $\mathbf{A}$ neither positive nor negative definite. This
invalidates the convergence guarantees of classical stationary
methods, which require diagonal dominance or definiteness
\cite{SaadBook2003}.

\paragraph{Spectral growth with frequency.}
As $\omega$ increases the diagonal shift $k^2 = \omega^2$ grows,
pulling eigenvalues further from the origin. The spectral radius
expands and the condition number $\kappa(\mathbf{A})$ deteriorates.
\todo{Add $\kappa(\mathbf{A})$ vs.\ $\omega$ from solver benchmarks.}
Krylov methods can still converge, but they require a large
subspace or a preconditioner that clusters eigenvalues near unity.
Constructing such a preconditioner for high-$\omega$ Helmholtz is
precisely the open problem that motivates this work: rather than
building a classical algebraic preconditioner, we ask whether a
neural operator trained on low-frequency solutions can serve as a
transferable one.

\paragraph{Non-normality.}
The PML introduces complex off-diagonal entries that break
self-adjointness. For non-normal matrices the spectrum alone
does not determine convergence: the pseudospectrum
--- the region of $\mathbb{C}$ where $(z\mathbf{I} - \mathbf{A})^{-1}$
is large in norm --- can extend far beyond the eigenvalues
\cite{Trefethen2005}. This makes spectral analysis insufficient
and motivates empirical convergence studies in
Chapter~\ref{chapter:classical_solvers}.

\begin{figure}[htbp]
    \centering
    \includegraphics[width=0.9\textwidth]{figures/c - eigenvalue scatter.png}
    \caption{Eigenvalue distribution of $\mathbf{A}(\omega)$ for
    $\omega \in \{16, 32, 64, 128\}$. The shaded region ($\operatorname{Re}(\lambda) < 0$)
    marks indefiniteness. As frequency increases, the spectrum shifts
    and spreads, violating the requirements for classical stationary
    iterative methods.}
    \label{fig:eigenvalues}
\end{figure}


% ============================================================
\section{Resolution requirements}
\label{sec:resolution_requirements}
% ============================================================

A finite difference solution is physically meaningful only if the
grid resolves the wavelength $\lambda = 2\pi c/\omega$. The
standard criterion requires at least $G$ points per wavelength
(PPW):
\begin{equation}
    h \;\leq\; \frac{2\pi c}{\omega G} = \frac{\lambda}{G}.
    \label{eq:ppw}
\end{equation}
For second-order stencils the accepted minimum is $G = 10$
\cite{Ihlenburg1998}; in practice a larger $G$ is advisable at
higher frequencies to limit the pollution error discussed in
Section~\ref{sec:dispersion_pollution}.

Throughout this work $n = 512$ is fixed. With $h = 1/(n-1)
\approx 1.96 \times 10^{-3}$, the four frequencies all exceed
the minimum criterion by a large margin
(Table~\ref{tab:resolution}). Fixing $n$ rather than $h$
deliberately isolates the effect of increasing $\omega$ on solver
behaviour: the system size stays constant while the spectral
difficulty grows. This is the experimental design choice that
makes the iteration-count comparisons in
Chapter~\ref{chapter:classical_solvers} interpretable.

\begin{table}[h]
\centering
\begin{tabular}{lllll}
\toprule
$\omega$ & $\lambda = 2\pi/\omega$ & $h$ &
PPW $= \lambda/h$ & Status \\
\midrule
16  & $0.393$ & $1.96 \times 10^{-3}$ & $\approx 200$ & Over-resolved \\
32  & $0.196$ & $1.96 \times 10^{-3}$ & $\approx 100$ & Well-resolved \\
64  & $0.098$ & $1.96 \times 10^{-3}$ & $\approx 50$  & Resolved \\
128 & $0.049$ & $1.96 \times 10^{-3}$ & $\approx 25$  & Marginally resolved \\
\bottomrule
\end{tabular}
\caption{Points per wavelength for the $n = 512$ grid with
$c = 1$. All four frequencies exceed the $G = 10$ minimum
criterion. The PPW decreases by a factor of eight from $\omega = 16$
to $\omega = 128$, driving the pollution error growth discussed
in Section~\ref{sec:dispersion_pollution}.}
\label{tab:resolution}
\end{table}

The PML occupies the outer $n_{\mathrm{pml}} = 112$ grid points
on each side. After masking, the usable interior domain is
$288 \times 288$ grid points. All loss functions and convergence
metrics in this work are evaluated on the interior only: the
attenuated PML solution carries no physical information and
polluting the metric with it would artificially lower the
reported error.


% ============================================================
\section{Dispersion and pollution errors}
\label{sec:dispersion_pollution}
% ============================================================

Satisfying \eqref{eq:ppw} controls the \emph{local} truncation
error but does not prevent global phase error from accumulating
over large domains. Two related phenomena matter.

\paragraph{Numerical dispersion.}
Substituting a plane wave $u_{i,j} = e^{i(p_x ih + p_y jh)}$
into \eqref{eq:fd_stencil} gives a modified dispersion relation
with discrete phase velocity $c_h \neq c$; the deviation is
$c_h - c = \mathcal{O}(h^2)$ at each point.

\paragraph{Pollution effect.}
Local accuracy does not imply global accuracy at high frequency.
Ihlenburg and Babu\v{s}ka \cite{Ihlenburg1995,Ihlenburg1998}
showed that the relative $H^1$ error satisfies
\begin{equation}
    \frac{\|u - u_h\|_{H^1}}{\|u\|_{H^1}}
    \;\leq\;
    C_1\,(kh)^2 \;+\; C_2\,k^3 h^2,
    \label{eq:pollution_bound}
\end{equation}
where $k = \omega/c$. The first term is the standard
interpolation error, controlled by keeping $kh$ small. The
second term is the \emph{pollution term}: for fixed PPW (i.e.\
$kh$ constant) it grows linearly with $k$. The mechanism is a
frequency-dependent deterioration of the inf-sup constant of
$\mathbf{A}$; the full argument is in \cite{Ihlenburg1998}.

Since $n = 512$ is fixed, the PPW decreases from approximately
200 at $\omega = 16$ to 25 at $\omega = 128$
(Table~\ref{tab:resolution}), so the pollution term $C_2 k^3 h^2$
grows with $\omega$. This is intentional: it ensures that
higher-frequency problems are genuinely harder to solve --- the
discrete system is not just larger, it is algebraically more
ill-conditioned --- which is precisely the scenario motivating
frequency transfer as a warm-start strategy.

\todo[inline]{Insert Fig.~3.4: log-log plot of relative
$H^1$ error versus $\omega$ at fixed $n = 512$, showing the
super-linear growth predicted by \eqref{eq:pollution_bound}.}


% ============================================================
\section{Absorbing boundary conditions}
\label{sec:pml_discretisation}
% ============================================================

The Helmholtz equation is posed on an unbounded domain. Hard
Dirichlet or Neumann conditions at the grid boundary produce
spurious reflections that corrupt the interior solution. Instead,
a Perfectly Matched Layer (PML) surrounds the domain: a boundary
strip in which outgoing waves are exponentially absorbed without
any reflection at the interior--layer interface
\cite{Berenger1994,ChewWeedon1994}.

\paragraph{Coordinate stretching.}
Inside the PML strip, the real coordinate $x$ is complexified via
the stretching function
\begin{equation}
    s(x) = 1 + i\,\sigma(x),
    \label{eq:stretching}
\end{equation}
where $\sigma(x) \geq 0$ vanishes in the interior and increases
into the strip. The physical effect is clean: a plane wave
$e^{ikx}$ travelling into the layer accumulates an imaginary
phase $\int\sigma\,\mathrm{d}x$, which means exponential decay
in amplitude without any reflection, because the medium is
matched to the interior. In the discrete operator, each
off-diagonal stencil weight at a PML grid point is divided
by $s(x_j)$:
\begin{equation}
    \frac{u_{i,j+1} - 2u_{i,j} + u_{i,j-1}}{h^2}
    \;\longrightarrow\;
    \frac{1}{s(x_j)\,h^2}
    \bigl(u_{i,j+1} - 2u_{i,j} + u_{i,j-1}\bigr).
\end{equation}
The sparsity pattern of $\mathbf{A}$ is unchanged; the entries
become complex-valued inside the PML region.

\paragraph{Damping profile.}
The damping grows quadratically from zero at the interior--PML
interface to the peak $\sigma_0$ at the outer wall,
\begin{equation}
    \sigma(x)
    = \sigma_0\!\left(\frac{d(x)}{n_{\mathrm{pml}}}\right)^{\!2},
    \label{eq:pml_profile}
\end{equation}
where $d(x)$ is the distance from the interior boundary in grid
points. With a Dirichlet wall at the outer edge, a plane wave
makes a round trip through the full layer, giving a reflection
coefficient
\begin{equation}
    R \;\approx\; e^{-\,2\sigma_0\,n_{\mathrm{pml}}\,h/3},
    \label{eq:R_formula}
\end{equation}
where the factor $1/3$ comes from integrating the quadratic
profile. Targeting $R = 10^{-3}$ sets the required product
$\sigma_0\,n_{\mathrm{pml}}\,h \geq \tfrac{3}{2}\ln(10^3) \approx 10.4$.

\paragraph{Choosing the layer thickness.}
The most natural design for a PML that must handle four
different frequencies would adapt \emph{both} the thickness
$n_{\mathrm{pml}}$ and the peak damping $\sigma_0$ to each
$\omega$. Shorter wavelengths need a physically larger absorbing
region to contain the same number of wavelengths; a systematic
search over a leakage criterion --- the ratio of wave amplitude
at the physical boundary to amplitude in the interior --- with
GMRES convergence as a secondary constraint found
frequency-optimal thicknesses of approximately
$\{104,\;84,\;68,\;56\}$ cells for $\omega \in \{16,32,64,128\}$.

This frequency-adaptive thickness was not adopted. The reason
goes to the heart of the learning problem. Varying
$n_{\mathrm{pml}}$ changes the total grid footprint and, with
it, the size of the physical interior domain. Each frequency
would then present the neural network with a different spatial
layout: the $288 \times 288$ interior at $\omega = 16$ would
become something else at $\omega = 128$. A transfer operator
trained on pairs $(u_{\omega_{\mathrm{low}}},\,u_{\omega_{\mathrm{high}}})$
from differently-sized grids would need to reason about
grid geometry as well as wave physics, conflating boundary
artefacts with genuine frequency dependence. The model would
risk memorising frequency-specific grid structure rather than
learning the transferable part of the solution map --- a
textbook path to overfitting.

Fixing
\begin{equation}
    n_{\mathrm{pml}} = 112 \quad (\text{all frequencies})
    \label{eq:fixed_npml}
\end{equation}
eliminates this confound. Every frequency has the same
$288 \times 288$ physical interior, the same sparsity pattern,
and the same input layout for the network. This choice also
makes the system size identical across frequencies, so that
iteration-count comparisons in
Chapter~\ref{chapter:classical_solvers} isolate the effect of
$\omega$ on solver convergence.

\paragraph{Frequency-adaptive peak damping.}
With $n_{\mathrm{pml}} = 112$ fixed, $\sigma_0$ is the only
remaining tuning knob. Two forces act in opposite directions.

\begin{itemize}
    \item \textit{Absorption.} Higher $\omega$ means more
    oscillations fit inside the fixed layer, so $\sigma_0$
    must increase with $\omega$ to maintain the target $R$.
    From \eqref{eq:R_formula}, keeping $R$ constant at fixed
    $n_{\mathrm{pml}}\,h$ requires $\sigma_0 \propto \omega$.

    \item \textit{Matrix conditioning.} A large $\sigma_0$
    inflates the imaginary part of the diagonal entries of
    $\mathbf{A}$, dragging the spectrum away from the real
    axis. This extends the pseudospectrum and increases GMRES
    iteration counts. In the extreme it can make the system
    harder to solve than without any PML at all.
\end{itemize}

The values
\begin{equation}
    \sigma_0 \;\in\; \{42.5,\;85.0,\;120.0,\;180.0\}
    \quad\text{for}\quad
    \omega \;\in\; \{16,\;32,\;64,\;128\}
    \label{eq:sigma0_values}
\end{equation}
were determined empirically by solving a grid of $(n_{\mathrm{pml}},
\sigma_0)$ configurations and measuring both the leakage metric
and GMRES convergence (Table~\ref{tab:system_params}). The
values grow with $\omega$ --- roughly doubling across the first
octave, then more slowly --- reflecting a deliberate sub-linear
trend that limits imaginary diagonal growth. No clean power law
fits all four points; the values represent a direct numerical
trade-off rather than a closed-form prediction, and are held
fixed across all experiments.

\begin{table}[h]
\centering
\begin{tabular}{lllll}
\toprule
$\omega$ & $\sigma_0$ & $n_{\mathrm{pml}}$ & Interior size & Target $R$ \\
\midrule
16  & $42.5$  & $112$ & $288^2$ & $10^{-3}$ \\
32  & $85.0$  & $112$ & $288^2$ & $10^{-3}$ \\
64  & $120.0$ & $112$ & $288^2$ & $10^{-3}$ \\
128 & $180.0$ & $112$ & $288^2$ & $10^{-3}$ \\
\bottomrule
\end{tabular}
\caption{PML parameters for the $n = 512$ grid. Fixing
$n_{\mathrm{pml}} = 112$ keeps the interior size identical
across all frequencies; only $\sigma_0$ is adapted. The sub-linear
growth of $\sigma_0$ with $\omega$ balances absorption quality
against matrix conditioning.}
\label{tab:system_params}
\end{table}

\begin{figure}[htbp]
    \centering
    \subfloat[1D cross-sections of $\sigma(x)$ for $\omega \in \{16, 32, 64, 128\}$.]{
        \includegraphics[width=0.8\textwidth]{figures/c - pml 1d.png}
        \label{fig:pml_1d}
    }
    \hfill
    \subfloat[2D spatial distribution of the quadratic damping on the $512^2$ grid.]{
        \includegraphics[width=0.8\textwidth]{figures/c - pml 2d.png}
        \label{fig:pml_2d}
    }
    \vspace{1em}
    \caption{PML configuration. (a) Quadratic ramp of $\sigma$ from
    zero at the interior interface (grid indices 112 and 400) to $\sigma_0$
    at the outer wall. Higher frequencies carry a larger $\sigma_0$ to
    maintain the target $R = 10^{-3}$. (b) The frame-like PML surrounding
    the $288 \times 288$ physical interior --- identical in footprint for
    all four frequencies.}
    \label{fig:pml_combined}
\end{figure}

\todo[inline]{Remake Fig.~\ref{fig:pml_combined} with a colour scale
and Arial font consistent with the rest of the thesis.}


% ============================================================
\section{Summary of the discrete system}
\label{sec:linear_system}
% ============================================================

At each frequency $\omega$ the construction above yields a sparse
complex linear system $\mathbf{A}(\omega)\mathbf{u} = \mathbf{f}$
with the following properties:

\begin{itemize}
    \item $N = 262{,}144$ unknowns, $\mathcal{O}(N)$ non-zeros.
    \item Indefinite, non-Hermitian spectrum whose condition number
          grows with $\omega$.
    \item Fixed PML geometry ($n_{\mathrm{pml}} = 112$) with
          frequency-adapted $\sigma_0$ \eqref{eq:sigma0_values}
          that keeps boundary reflections below $10^{-3}$
          without over-stiffening the matrix.
    \item A $288 \times 288$ physical interior identical across
          all frequencies; the PML region is excluded from all
          error metrics and training losses.
\end{itemize}

The system is deliberately designed so that increasing $\omega$
increases algebraic difficulty without changing system size or
spatial layout. This makes the question concrete: \emph{can a
neural operator trained at low frequency provide a useful warm
start for the iterative solver at high frequency, reducing the
iteration count?} Chapter~\ref{chapter:classical_solvers}
benchmarks state-of-the-art classical solvers on this system
and shows that the answer to that question cannot come from
classical preconditioners alone.

\begin{figure}[htbp]
    \centering
    \includegraphics[width=0.85\textwidth]{figures/c - condition.png}
    \caption{Condition number $\kappa(\mathbf{A}(\omega))$ versus
    $\omega$ for constant $\sigma_0$ (dashed) and adaptive $\sigma_0$
    \eqref{eq:sigma0_values} (solid). Adaptive damping homogenises the
    algebraic difficulty across frequencies, so that differences in
    solver performance reflect physical frequency dependence rather
    than PML artefacts.}
    \label{fig:condition_homogenisation}
\end{figure}
